\documentclass[literature.tex]{subfiles}

\begin{document}
\prose{Šifra Mistra Leonarda}{Dan Brown}
\teorieA{
epická próza, dobrodružný román, konspirační thriller s detektivními prvky
}{
Populární (konzumní) literatura navazující na postmodernismus. Kombinuje historická fakta s fikcí a konspiračními teoriemi.
}{
Dynamický, čtivý a akční text s krátkými větami pro udržení tempa. Bohatá odborná terminologie z oblasti umění, historie, náboženství a kryptologie (francouzština, latina, staroangličtina). Časté cliffhangery a zvraty.
}{
Opakování motivů, ironie, cliffhanger, dialogy s gradací.
}{
Konspirace, symbolismus (pentagram, Mona Lisa, Vitruviova figura), alegorie (Svatý grál), kontrast mezi oficiální historií a „skrytou pravdou“.
}{
Rychlý, filmový styl s krátkými kapitolami a střídáním akce, hádanek a vysvětlujících pasáží. Autor často používá cliffhangery na konci kapitol pro udržení napětí.
}{
Langdon nemohl odtrhnout oči od fotografie, kterou držel v ruce, a pocítil neurčitý strach.
Obrázek byl příšerný a hluboce podivný, vyvolával zneklidňující dojem jakéhosi \wemph{déjá-vu}.
Policista se podíval na hodinky.
\enquote{Můj kapitán vás očekává, pane.}
Langdon ho sotva slyšel.
Očima se vpíjel do obrázku.
\enquote{Ten symbol tady, a~to, jak~podivně je~to~tělo…} \texthl{\enquote{Naaranžováno?}} pomohl mu policista.
Langdon přikývl a~projel jím~chlad.
\enquote{Neumím~si~představit, kdo může tohle udělat jinému člověku.}
Policista vypadal zachmuřeně.
\enquote{Vy to nechápete, pane Langdone. To,~co~vidíte na té fotografii…} Odmlčel se.
\wemph{\enquote{Monsieur~Sauniére si~to~udělal~sám.}}
}
\teorieB{
\wtem{Robert Langdon}{profesor symbolologie z Harvardu, racionální, skeptický, ale rychle zatažen do víru událostí.}
\wtem{Sophie Neveuová}{kryptoložka pařížské policie, vnučka Jacquese Saunièra, chytrá a odvážná.}
\wtem{Sir Leigh Teabing (Učitel)}{britský historik a znalec grálu, zdánlivě pomocník, ve skutečnosti hlavní manipulátor.}
\wtem{Silas}{albín, mnich Opus Dei, tragická postava fanaticky oddaná víře.}
\wtem{Bezu Fache}{kapitán pařížské policie, pragmatický vyšetřovatel.}
\wtem{Jacques Sauni\`{e}re}{kurátor Louvru, velmistr Převorství sionského, dědeček Sophie.}
}{
V Louvru je zavražděn kurátor Jacques Saunière. Na místě činu je nalezena série šifer a symbolů. Policie povolá symbologa Roberta Langdona, který je brzy obviněn z vraždy. Pomáhá mu kryptoložka Sophie Neveuová (Saunièrova vnučka). Společně luští hádanky odkazující na díla Leonarda da Vinciho, Převorství sionské a tajemství Svatého grálu. Na stopu se jim dostanou fanatičtí členové Opus Dei (Silas) i záhadný Učitel. Pátrání je zavede z Paříže do Londýna a nakonec odhalí kontroverzní interpretaci křesťanských dějin.
}{
Rychlý, filmový styl s krátkými kapitolami a častými cliffhangery. Prostor: Paříž (Louvre), Londýn (Temple Church, Rosslyn). Čas: současnost (začátek 21. století), s častými historickými odbočkami.
}{
Kombinace napětí a konspirační teorie. Kniha zpochybňuje oficiální dějiny křesťanství (úloha Máří Magdaleny, vznik Nového zákona). Hlavní myšlenka: pravda je často skrytá za oficiální verzí dějin a mocenské instituce (církev) ji dlouhodobě potlačují. Zároveň oslava intelektu, hádanek a hledání pravdy.
}
\historie{
Kniha vyšla v roce 2003 v období silného zájmu o konspirační teorie (po 11. září) a kritiky institucionálního náboženství.
}{
Inspirace historickými konspiračními teoriemi (Převorství sionské, Opus Dei, Svatý grál), díly Leonarda da Vinciho a knihami jako \textit{Svatá krev a Svatý grál}.
}{}{}{}{
Dan Brown (*1964). Americký spisovatel. Původně učitel angličtiny a zpěvák. Proslavil se sérií thrillerů s Robertem Langdonem. Da Vinciho kód se stal celosvětovým bestsellerem (přes 80 milionů prodaných výtisků).
}{
Mezi nejznámější díla patří \textsc{Andělé a démoni}, \textsc{Iluze}, \textsc{Inferno}, \textsc{Původ}.
}\newpage
\kritika{
Obrovský komerční úspěch, ale silná kritika ze strany katolické církve a historiků za nepřesnosti a „vymýšlení faktů“. Kniha vyvolala celosvětovou debatu o vztahu fikce a reality.
}{}{
Konspirační thrillery zůstávají populární. Téma zpochybňování oficiálních narativů (církev, dějiny) rezonuje i dnes v éře fake news a konspiračních teorií.
}
\subwsection{Kontrola faktů:}
\begin{indentpara}
\textbl{Opus Dei} -- skutečná katolická organizace, známá konzervativním přístupem, ale Brown ji silně démonizuje.\\
\textbl{Převorství sionské} -- moderní podvod (vymyšleno v 50. letech 20. století), Brown jej prezentuje jako středověký řád.
\end{indentpara}
\nazor{
Da Vinciho kód je perfektní „page-turner“ – čte se jedním dechem díky krátkým kapitolám a neustálým zvratům. Nejvíce mě baví hádanky spojené s uměním Leonarda da Vinciho a rychlé tempo. Na druhou stranu je zřejmé, že Brown hodně „upravuje“ historii, což může u méně informovaných čtenářů vést k zavádějícím představám. Přesto jako čistě zábavný konspirační thriller funguje výborně. Kniha je nenáročná, napínavá a originální – ideální oddechovka, která dokáže člověka vtáhnout na několik večerů.
}
\explain{
\textbl{Dobrodružný román --} žánr zaměřený na napínavý děj a hrdinské činy;
\textbl{Konspirační thriller --} kombinuje napětí s odhalováním tajných spiknutí;
\textbl{Populární literatura --} masově přístupná, komerčně úspěšná beletrie.
}
\wpage
\end{document}